\chapter{The Least-Squares Core: Givens Rotations and the Running QR}
\label{ch:givens}

% local vector macros for the running-QR right-hand side (g) and Hessenberg column (h)
\providecommand{\vg}{\mathbf{g}}
\providecommand{\vh}{\mathbf{h}}

\chapterorient{GMRES, the engine of \cref{ch:gmres}, leaves a small chore at the
end of every step: minimize $\norm{\beta\,\vec e_1 - \bar{\mat H}_m\vy}$ over a tiny
least-squares problem whose matrix grows by one column each iteration. This chapter
is about that chore---and about the surprisingly elegant way the solver does it. We
do not re-solve the least-squares problem from scratch at each step; we maintain a
\emph{running} QR factorization, updated one column at a time by a stream of $2\times2$
plane rotations. The stream is a fold over the arriving columns, in the exact sense of
\cref{ch:combinators}. And it has a gift: after the rotation hits the right-hand side,
the GMRES residual norm is sitting there, already computed, for free. No
$\mat A\vx_m$, no residual evaluation---a unitary byproduct of triangularizing the
matrix. By the end you will be able to generate a Givens rotation by hand, run the
stream on a small Hessenberg, read the residual off one entry, and back-solve for the
answer: the entire inner loop of GMRES, made concrete.}

\section{The chore GMRES leaves behind}
\label{sec:gv-setup}

Recall the shape of GMRES from \cref{ch:gmres}. At step $m$ the Arnoldi process of
\cref{ch:orthog} has built an orthonormal basis $\mat V_m = [\vv_1,\dots,\vv_m]$ for
the Krylov subspace and, as a byproduct of orthogonalizing each new direction, an
$(m{+}1)\times m$ matrix $\bar{\mat H}_m$ that records the orthogonalization
coefficients. This matrix is \emph{upper-Hessenberg}: upper-triangular except for a
single nonzero sub-diagonal, the band of entries $h_{k+1,k}$ just below the main
diagonal. GMRES then forms its iterate as $\vx_m = \vx_0 + \mat V_m\vy$, choosing the
coordinate vector $\vy\in\Real^m$ that minimizes the residual. Because Arnoldi makes
the residual norm equal to a small projected quantity, that minimization collapses to
\begin{equation}
\label{eq:gv-lsq}
  \min_{\vy\in\Real^m}\ \norm{\beta\,\vec e_1 - \bar{\mat H}_m\,\vy},
\end{equation}
where $\beta = \norm{\vr_0}$ is the initial residual norm and $\vec e_1$ is the first
standard basis vector. This is a \emph{least-squares} problem---more equations
($m{+}1$) than unknowns ($m$)---and it is \emph{small}: its size is the iteration
count, never the dimension $N$ of the field. A GMRES run that converges in 40 steps
solves a $41\times 40$ least-squares problem, whatever the millions of unknowns in
$\vx$.

\Cref{eq:gv-lsq} is the chore. We could, at every step, assemble the current
$\bar{\mat H}_m$, factor it afresh, and solve. That is correct and wasteful: the
matrix at step $m{+}1$ is the matrix at step $m$ with one column glued on the right
and one row glued on the bottom; almost all the work of factoring it was already done
last step. The professional move is to maintain a factorization \emph{incrementally},
spending only the cost of the one new column each step. That incremental
factorization is the subject of this chapter, and the tool that builds it is the
plane rotation.

\begin{figure}[t]
\centering
\begin{tikzpicture}[font=\footnotesize]
  % the Hessenberg with one column arriving and one sub-diagonal to kill
  \def\cw{0.62} \def\rh{0.52}
  % grid 5 rows x 4 cols
  \foreach \c in {0,1,2,3}{
    \foreach \r in {0,...,4}{
      \draw[simGray!35] (\c*\cw,-\r*\rh) rectangle ++(\cw,-\rh);
    }
  }
  % triangular block already done (cols 0..2): shade upper-tri
  \foreach \c/\maxr in {0/1,1/2,2/3}{
    \foreach \r in {0,...,\maxr}{
      \fill[simBgBlue] (\c*\cw,-\r*\rh) rectangle ++(\cw,-\rh);
    }
  }
  % the arriving column (col 3): shade gold, full height down to its sub-diagonal r=4
  \foreach \r in {0,...,4}{
    \fill[simBgGold] (3*\cw,-\r*\rh) rectangle ++(\cw,-\rh);
  }
  % redraw gridlines over fills
  \foreach \c in {0,1,2,3}{\foreach \r in {0,...,4}{
      \draw[simGray!35] (\c*\cw,-\r*\rh) rectangle ++(\cw,-\rh);}}
  % mark the sub-diagonal entry to be annihilated in the new column: (row 4, col 3)
  \draw[simRust, very thick] (3*\cw,-4*\rh) rectangle ++(\cw,-\rh);
  \node[font=\scriptsize, text=simRust] at (3*\cw+0.5*\cw,-4*\rh-0.5*\rh) {$h_{j+1,j}$};
  \node[font=\scriptsize] at (3*\cw+0.5*\cw,-3*\rh-0.5*\rh) {$h_{j,j}$};
  % bracket labels
  \node[font=\scriptsize, text=simBlue, align=center] at (1*\cw,0.55) {already triangular\\(cols $0..j{-}1$)};
  \node[font=\scriptsize, text=simGold!70!black, align=center] at (3*\cw+0.5*\cw,0.55) {new\\column $j$};
  % the rotation arrow: zero the rust cell against the one above it
  \draw[-{Stealth[length=2mm]}, simRust, thick]
    (3*\cw+\cw+0.15,-4*\rh-0.5*\rh) .. controls +(0.7,0.2) and +(0.7,-0.2) ..
    (3*\cw+\cw+0.15,-3*\rh-0.5*\rh);
  \node[font=\scriptsize, text=simRust, anchor=west] at (3*\cw+\cw+0.55,-3.5*\rh-0.5*\rh)
    {$G_j$ rotates};
  \node[font=\scriptsize, text=simRust, anchor=west] at (3*\cw+\cw+0.55,-4.0*\rh-0.5*\rh)
    {$h_{j+1,j}\!\to\!0$};
\end{tikzpicture}
\caption[A Givens rotation zeroing one sub-diagonal entry]{The geometry of one step.
The upper-Hessenberg matrix $\bar{\mat H}_m$ is triangular in its leading columns
(blue), with a single nonzero sub-diagonal. A new column $j$ arrives from Arnoldi
(gold) carrying one fresh sub-diagonal entry $h_{j+1,j}$ (rust outline). A single
$2\times2$ plane rotation $G_j$, mixing rows $j$ and $j{+}1$, annihilates that entry
against the diagonal above it---completing the triangularization of column $j$. One
rotation, one sub-diagonal entry: this is the atom the rest of the chapter assembles
into a stream.}
\label{fig:gv-hessenberg}
\end{figure}

\section{Givens rotations: zeroing one entry at a time}
\label{sec:gv-rotations}

A \emph{Givens rotation}, or plane rotation, is a $2\times2$ orthogonal matrix
\begin{equation}
\label{eq:gv-matrix}
  G(c,s) = \begin{pmatrix} c & s \\ -s & c \end{pmatrix},
  \qquad c^2 + s^2 = 1,
\end{equation}
that rotates the plane by the angle whose cosine is $c$ and sine is $s$. Its single
job in this chapter is surgical: given a 2-vector $\begin{psmallmatrix} a \\ b
\end{psmallmatrix}$, choose $c$ and $s$ so that applying $G$ sends the second
component to zero,
\begin{equation}
\label{eq:gv-goal}
  \begin{pmatrix} c & s \\ -s & c \end{pmatrix}
  \begin{pmatrix} a \\ b \end{pmatrix}
  = \begin{pmatrix} r \\ 0 \end{pmatrix}.
\end{equation}
Because $G$ is orthogonal it preserves length, so the surviving first component must
carry the whole norm: $r = \sqrt{a^2 + b^2}$. That single observation \emph{is} the
construction. The bottom row of \cref{eq:gv-goal} reads $-s\,a + c\,b = 0$, and
together with $r = \sqrt{a^2+b^2}$ the top row reads $c\,a + s\,b = r$. Solving the
pair gives the rotation coefficients directly.

\begin{lemma}[Givens generation]
\label{lem:gv-generate}
For a real 2-vector $(a,b)$ with $r = \sqrt{a^2 + b^2} > 0$, the rotation $G(c,s)$
with
\begin{equation}
\label{eq:gv-cs}
  c = \frac{a}{r}, \qquad s = \frac{b}{r}
\end{equation}
satisfies $G(c,s)\begin{psmallmatrix} a \\ b \end{psmallmatrix} =
\begin{psmallmatrix} r \\ 0 \end{psmallmatrix}$, and $c^2 + s^2 = 1$.
\end{lemma}

\begin{proof}
The orthogonality identity is immediate: $c^2 + s^2 = (a^2 + b^2)/r^2 = 1$. For the
action, compute both rows. The top row gives $c\,a + s\,b = (a^2 + b^2)/r = r^2/r = r$.
The bottom row gives $-s\,a + c\,b = (-b\,a + a\,b)/r = 0$. Hence the image is exactly
$(r,0)^{\!\top}$.
\end{proof}

We call the construction of $(c,s)$ from $(a,b)$ the verb \code{givens\_generate}, and
the in-place action of a stored $(c,s)$ on a pair of scalars the verb
\code{givens\_apply}. The two are the entire kernel; everything later in the chapter is
a way of \emph{scheduling} calls to them.

\begin{lstlisting}[style=l4style]
givens_generate :: (a: Scalar, b: Scalar) -> (c: Real, s: Scalar)
givens_apply    :: (c: Real, s: Scalar) -> ((dx: Scalar, dy: Scalar) -> (Scalar, Scalar))
givens_apply c s (dx, dy) = (c*dx + s*dy, -conj s * dx + c*dy)
\end{lstlisting}

The action of \code{givens\_apply} is the general orthogonal update $\begin{psmallmatrix}
dx' \\ dy' \end{psmallmatrix} = \begin{psmallmatrix} c & s \\ -\bar s & c
\end{psmallmatrix}\begin{psmallmatrix} dx \\ dy \end{psmallmatrix}$, which over the
reals is \cref{eq:gv-matrix} exactly; the conjugate $\bar s$ appears only in the
complex case, of which more below. Note the \emph{closure-returning} signature of
\code{givens\_apply}: it takes the rotation $(c,s)$ and \emph{returns a function} of the
pair, written with the paren-grouping of \cref{ch:combinators}. This is deliberate. A
generated rotation is applied many times---to its own column, to later columns, to the
right-hand side---so it is natural to read \code{givens\_apply}~$c$~$s$ as ``the
operator that rotates a pair,'' a first-class value stored and reused.

\begin{notationbox}
Throughout this chapter $(c,s)$ is one stored rotation; $a$ and $b$ are the two
components of the 2-vector it is built to annihilate (here $a$ is the diagonal, $b$ the
sub-diagonal entry to be zeroed); $r = \sqrt{a^2+b^2}$ is the surviving norm. We index
the rotation generated at Arnoldi step $j$ as $(c_j, s_j)$, and we write a generic
applied pair as $(dx, dy)$. The cosine $c$ is always real; the sine $s$ is real for a
real system and complex for a complex one.
\end{notationbox}

\subsection{The stable formula}
\label{sec:gv-stable}

\Cref{eq:gv-cs} is correct but dangerous to compute literally. Forming $r =
\sqrt{a^2 + b^2}$ squares each component, and if $a$ or $b$ is larger than the square
root of the floating-point overflow threshold---around $10^{154}$ in double
precision---the intermediate $a^2$ overflows to infinity even though $r$ itself is a
perfectly representable number. The symmetric failure happens for tiny values, which
underflow to zero. Production libraries therefore never square before scaling. They
factor out the larger component first:
\begin{equation}
\label{eq:gv-scaled}
  t = \frac{\min(\abs a,\abs b)}{\max(\abs a,\abs b)}, \qquad
  r = \max(\abs a,\abs b)\,\sqrt{1 + t^2},
\end{equation}
so the square root sees an argument between $1$ and $2$ and the division $t$ is between
$0$ and $1$. Mathematically this is identical to $r = \sqrt{a^2+b^2}$; numerically it
never overflows unless the true answer does. Palace follows the LAPACK convention
(the algorithm of Bindel, Demmel, Kahan, and Marques), computing $c = \abs a / r$ and
fixing the sign of $r$ to match $a$ so that the diagonal stays positive. We treat the
scaling as a \emph{load-bearing numerical detail}, not a transparent rewrite of
\cref{eq:gv-cs}: it changes which inputs the algorithm survives, and the precise scaled
form is recorded in \cref{app:numerics}.

\begin{pitfall}
Do not implement \cref{eq:gv-cs} as written. The literal formula
$r=\sqrt{a^2+b^2}$ overflows for large inputs and underflows for small ones, returning
$\infty$ or $0$ where the scaled form \eqref{eq:gv-scaled} returns the right finite
$r$. This is not a micro-optimization; it is a correctness fix at the edges of the
floating-point range, and it is exactly the kind of ``load-bearing numerical trick''
\cref{ch:bigidea} warned must be preserved as part of the algorithm rather than
abstracted away.
\end{pitfall}

\begin{workedexample}{Generating a Givens rotation by hand}{gv-generate}
Build the rotation that annihilates the second component of $(a,b) = (3,4)$, and check
its action.

\emph{The norm.} $r = \sqrt{3^2 + 4^2} = \sqrt{9 + 16} = \sqrt{25} = 5$. (The
$(3,4,5)$ triangle keeps the arithmetic clean.)

\emph{The coefficients.} By \cref{eq:gv-cs},
\[
  c = \frac{a}{r} = \frac{3}{5} = 0.6, \qquad
  s = \frac{b}{r} = \frac{4}{5} = 0.8.
\]
Sanity check the orthogonality identity: $c^2 + s^2 = 0.36 + 0.64 = 1$. Good---this is
a genuine rotation.

\emph{The action.} Apply $G(0.6,0.8)$ to $(3,4)$ using \code{givens\_apply}:
\[
  \begin{pmatrix} 0.6 & 0.8 \\ -0.8 & 0.6 \end{pmatrix}
  \begin{pmatrix} 3 \\ 4 \end{pmatrix}
  = \begin{pmatrix} 0.6\cdot 3 + 0.8\cdot 4 \\ -0.8\cdot 3 + 0.6\cdot 4 \end{pmatrix}
  = \begin{pmatrix} 1.8 + 3.2 \\ -2.4 + 2.4 \end{pmatrix}
  = \begin{pmatrix} 5 \\ 0 \end{pmatrix}.
\]
The second component is exactly zero and the first is exactly $r = 5$: the whole length
of the original 2-vector survived in the surviving component, as orthogonality
demands. The scaled formula \eqref{eq:gv-scaled} would compute the same $(c,s)=(0.6,0.8)$
via $\max=4$, $t = 3/4$, $r = 4\sqrt{1 + 9/16} = 4\cdot(5/4) = 5$---identical answer,
no squaring of large numbers.
\end{workedexample}

\subsection{The complex case, briefly}
\label{sec:gv-complex}

For a complex linear system the Hessenberg entries are complex, and the rotation must
preserve the Hermitian (conjugate) norm. The convention Palace and LAPACK use keeps the
cosine $c$ \emph{real} and lets the sine $s$ be complex, with the unitarity identity
becoming $c^2 + \abs s^2 = 1$. The action of \code{givens\_apply} carries the conjugate
in the lower row, $dy' = -\bar s\,dx + c\,dy$, which is why we wrote it that way above.
Nothing structural changes: the generation still produces a $(c,s)$ that zeros the
sub-diagonal, the stream still triangularizes, and the residual still reads off the
right-hand side. The real-cosine convention is a tidy bookkeeping choice---the rotation
register $c$ is one real number per step, $s$ one complex number---and we will not dwell
on it. The driven and eigenmode analyses (\cref{part:modalities}) supply the complex
systems that exercise it.

\section{The running QR stream}
\label{sec:gv-stream}

Now we schedule the rotations. The factorization we maintain is $\bar{\mat H}_m =
\mat Q_m \bar{\mat R}_m$, where $\mat Q_m$ is the accumulated product of all the plane
rotations and $\bar{\mat R}_m$ is upper-triangular. We never form $\mat Q_m$ as a
matrix; we store only the list of rotation pairs $(c_0,s_0),(c_1,s_1),\dots$, which is
all of $O(m)$ scalars. Each Arnoldi step delivers one new Hessenberg column, and the
stream processes it in four sub-steps.

\begin{description}[style=nextline, leftmargin=0pt, font=\normalfont\itshape]
  \item[(i) Replay.] The new column $\vh = \bar{\mat H}_m[:,j]$ arrives in the
  \emph{original} (un-rotated) frame. The previous columns were triangularized by the
  stored rotations $(c_0,s_0),\dots,(c_{j-1},s_{j-1})$; to bring the new column into the
  same frame we replay all of them, in order, against its entries:
  for $k=0,1,\dots,j{-}1$, apply $(c_k,s_k)$ to the pair $(h_k,h_{k+1})$.
  \item[(ii) Generate.] After replay, the entry $h_{j+1}$ on the sub-diagonal is the
  one fresh nonzero below the diagonal. Generate the new rotation $(c_j,s_j) =
  \code{givens\_generate}(h_j, h_{j+1})$ to annihilate it (\cref{lem:gv-generate}).
  \item[(iii) Apply to the column.] Apply the freshly generated $(c_j,s_j)$ to its own
  pair $(h_j, h_{j+1})$. This sets $h_{j+1} \gets 0$ and $h_j \gets r$: column $j$ is
  now upper-triangular, and $\bar{\mat R}_m$ has grown by a column.
  \item[(iv) Apply to the right-hand side.] Apply the \emph{same} $(c_j,s_j)$ to the
  right-hand-side pair $(g_j, g_{j+1})$, where $\vg$ began life as $\beta\,\vec e_1$.
  This is the step that pays the dividend, and \cref{sec:gv-residual} collects it.
\end{description}

\begin{figure}[t]
\centering
\begin{tikzpicture}[font=\footnotesize, node distance=5mm]
  \node[data, minimum width=18mm] (col) at (-0.2,0) {new column\\$\vh = \bar{\mat H}[:,j]$};
  \node[opbox, minimum width=20mm, right=8mm of col] (replay)
    {(i) \textbf{replay}\\\scriptsize $(c_0,s_0)\dots(c_{j-1},s_{j-1})$};
  \node[opbox, minimum width=20mm, right=7mm of replay, fill=simBgGold, draw=simGold] (gen)
    {(ii) \textbf{generate}\\\scriptsize $(c_j,s_j)$};
  \node[opbox, minimum width=18mm, below=9mm of gen] (apcol)
    {(iii) \textbf{apply}\\\scriptsize to column $\to \bar{\mat R}$};
  \node[opbox, minimum width=18mm, left=10mm of apcol] (aprhs)
    {(iv) \textbf{apply-rhs}\\\scriptsize to $\vg=\beta\vec e_1$};
  \node[product, minimum width=16mm, left=9mm of aprhs] (res)
    {residual\\$|g_{j+1}|$};
  \draw[flow] (col) -- (replay);
  \draw[flow] (replay) -- (gen);
  \draw[flow] (gen.south) -- (apcol.north);
  \draw[flow] (gen.south) .. controls +(0,-0.9) and +(0,0.9) .. (aprhs.north);
  \draw[flow] (aprhs) -- (res);
  % store-and-grow note for the new rotation
  \node[font=\scriptsize\itshape, text=simGreen, align=center] at (gen.north)[yshift=6mm]
    {append $(c_j,s_j)$\\to the stored list};
  \draw[refedge, simGreen] (gen.north) -- ++(0,4mm);
\end{tikzpicture}
\caption[The running-QR stream as a four-step pipeline]{The four sub-steps the stream
runs on each arriving Hessenberg column. \emph{Replay} (i) applies all stored rotations
to the new column, bringing it into the triangularized frame; \emph{generate} (ii)
produces one new rotation to kill the fresh sub-diagonal entry and appends it to the
stored list; \emph{apply} (iii) uses it to triangularize the column, growing
$\bar{\mat R}$; \emph{apply-rhs} (iv) uses the same rotation on the right-hand side
$\vg$, and the new last entry $|g_{j+1}|$ is the residual norm (\cref{sec:gv-residual}).
The pipeline runs once per Arnoldi step and the stored list grows by one rotation each
time.}
\label{fig:gv-pipeline}
\end{figure}

\begin{lstlisting}[style=l4style]
-- one column of the running QR (Arnoldi step j); see Figure (the four-step pipeline)
ls_update_column :: LsqState -> (h: HessCol) -> LsqState
ls_update_column st h =
  let h1       = foldl replayOne h (zip st.cs st.sn)      -- (i) replay stored 0..j-1
      (cj, sj) = givens_generate (h1 !! j) (h1 !! (j+1))  -- (ii) generate new rotation
      h2       = givens_apply cj sj (h1 !! j, h1 !! (j+1)) -- (iii) triangularize column j
      g'       = applyRhs cj sj st.g                       -- (iv) rotate the RHS pair
  in st { R = st.R `appendCol` h2          -- R grows by a column
        , cs = st.cs ++ [cj], sn = st.sn ++ [sj]  -- rotation list grows by one
        , g  = g'                          -- residual now lives in g'[j+1]
        }
  where replayOne col (ck, sk) = ...   -- givens_apply on pair (k, k+1)
\end{lstlisting}

\subsection{The stream is a fold over columns}
\label{sec:gv-fold}

Look at the signature of \code{ls\_update\_column}: it takes a state and a column and
returns a new state. That is exactly the shape of a fold step from \cref{ch:combinators}.
The whole running QR is
\begin{equation}
\label{eq:gv-fold}
\begin{aligned}
  \code{running\_qr}\ st_0\ [\vh_0, \dots, \vh_{m-1}]
  \;=\;\ &\code{foldl}\ \code{ls\_update\_column}\ st_0\\
  &[\vh_0, \vh_1, \dots, \vh_{m-1}],
\end{aligned}
\end{equation}
a left fold of \code{ls\_update\_column} over the stream of Hessenberg columns, with the
carry $st$ holding the triangular factor $\bar{\mat R}$, the growing rotation list
$(\code{cs},\code{sn})$, and the rotated right-hand side $\vg$. This is not a metaphor.
The carry threads strictly forward---column $j$ cannot be processed before column
$j{-}1$, because the replay sub-step (i) needs all the rotations the earlier columns
generated---so by the analysis of \cref{ch:combinators} the stream is an
\emph{intrinsically sequential} fold, not a parallel map. The rotation list is the part
of the carry that grows: one $(c_j,s_j)$ appended per step, so after $m$ columns the
carry holds $m$ rotations (\cref{fig:gv-growing}).

\begin{figure}[t]
\centering
\begin{tikzpicture}[font=\footnotesize]
  % three successive states of the carry, rotation list growing
  \foreach \step/\x/\n in {0/0/1, 1/4.2/2, 2/8.4/3}{
    \node[font=\scriptsize\bfseries, text=simBlue] at (\x+0.6,2.0) {after column $\step$};
    % stored rotation list as stacked cells
    \foreach \i in {0,...,\numexpr\n-1\relax}{
      \node[opbox, minimum width=14mm, minimum height=5mm] at (\x+0.6,1.4-0.62*\i)
        {$(c_{\i},s_{\i})$};
    }
  }
  % arrows between states: one column arrives, one rotation appended
  \draw[flow] (1.7,0.9) -- node[above,font=\scriptsize]{$\vh_1$} node[below,font=\scriptsize,text=simGreen]{$+\,(c_1,s_1)$} (3.0,0.9);
  \draw[flow] (5.9,0.9) -- node[above,font=\scriptsize]{$\vh_2$} node[below,font=\scriptsize,text=simGreen]{$+\,(c_2,s_2)$} (7.2,0.9);
  \node[font=\scriptsize\itshape, text=simGray, align=center] at (4.8,-1.2)
    {each Arnoldi column appends exactly one rotation to the carry:\\the fold's accumulator
     grows by one cell per step};
\end{tikzpicture}
\caption[The rotation list grows per Arnoldi column]{The stored rotation list is the
growing part of the fold's carry. Each Arnoldi column drives one
\code{ls\_update\_column} step, which appends exactly one new pair $(c_j,s_j)$ to the
list. After processing $m$ columns the carry holds $m$ rotations---an $O(m)$ store, far
cheaper than the explicit $\mat Q_m$ it stands in for. Because step $j$ replays all
earlier rotations, the list \emph{must} be threaded forward: this is a fold, not a map.}
\label{fig:gv-growing}
\end{figure}

\begin{pitfall}
The replay must come \emph{before} the generate, and the order of the replayed
rotations matters: $(c_0,s_0)$ then $(c_1,s_1)$ then $\dots$, never permuted. Plane
rotations do not commute---each one is built in the frame the previous ones
established---so applying them out of order, or generating the new rotation against an
un-replayed column, triangularizes against the wrong diagonal and silently corrupts
$\bar{\mat R}$. This is the running QR's load-bearing invariant: \emph{replay all stored
rotations in order, then generate.} It is the reason the carry is threaded forward and
the reason the stream is a fold rather than a map.
\end{pitfall}

\section{The residual, for free}
\label{sec:gv-residual}

Here is the payoff that makes the whole construction worth building. After sub-step
(iv)---applying the new rotation to the right-hand side $\vg$---the magnitude of the
\emph{new last entry} of $\vg$ is exactly the GMRES residual norm:
\begin{equation}
\label{eq:gv-free}
  \norm{\beta\,\vec e_1 - \bar{\mat H}_m\,\vy_m}
  \;=\; \abs{g_{j+1}},
\end{equation}
where $\vy_m$ is the (not-yet-computed) least-squares minimizer. We get the residual
without forming $\vy_m$, without forming the iterate $\vx_m = \vx_0 + \mat V_m\vy_m$,
and---most importantly---without the expensive matrix--vector product $\mat A\vx_m$ that
a literal residual evaluation would demand. The convergence test of GMRES is therefore a
single absolute value of a scalar the stream already computed.

Why is \cref{eq:gv-free} true? Because every rotation is orthogonal, and orthogonal maps
preserve length. Let $\mat Q_m$ be the accumulated rotation product, so that $\mat Q_m
\bar{\mat H}_m = \bar{\mat R}_m$ is upper-triangular and $\vg = \mat Q_m(\beta\,\vec
e_1)$. The least-squares residual is invariant under the orthogonal $\mat Q_m$:
\begin{equation}
\label{eq:gv-residual-derivation}
  \norm{\beta\,\vec e_1 - \bar{\mat H}_m\vy}
  = \norm{\mat Q_m\!\left(\beta\,\vec e_1 - \bar{\mat H}_m\vy\right)}
  = \norm{\vg - \bar{\mat R}_m\vy}.
\end{equation}
Now $\bar{\mat R}_m$ is $(m{+}1)\times m$ and upper-triangular: its last row is entirely
zero. So whatever $\vy$ we pick, the bottom equation $0 = g_{j+1} - 0$ can never be
satisfied, and the best we can do is zero out the top $m$ rows (which a nonsingular
triangular solve does exactly), leaving the residual equal to the one unkillable
entry $\abs{g_{j+1}}$. Splitting the norm into the part a triangular solve can zero and
the part it cannot:
\begin{equation}
\label{eq:gv-split}
  \norm{\vg - \bar{\mat R}_m\vy}^2
  = \underbrace{\norm{\vg_{0:m} - \mat R_m\vy}^2}_{\text{drive to }0\text{ by back-solve}}
  + \underbrace{\abs{g_{j+1}}^2}_{\text{the residual}}.
\end{equation}
The minimizer kills the first term, and the second term---the magnitude of the rotated
right-hand side's tail---is the residual. That is \cref{eq:gv-free}. The rotation stream,
built to triangularize the matrix, drags the right-hand side along and concentrates the
entire un-fittable residual into one entry, where we can read it.

\begin{figure}[t]
\centering
\begin{tikzpicture}[font=\footnotesize]
  % the rotated RHS vector g, last entry highlighted as the residual
  \def\rh{0.5}
  \node[font=\scriptsize\bfseries, text=simBlue] at (-1.6,0.3) {$\vg = \mat Q_m(\beta\vec e_1)$};
  \foreach \i/\lbl in {0/{g_0},1/{g_1},2/{\vdots},3/{g_m}}{
    \draw[simGray!45] (0,-\i*\rh) rectangle ++(0.9,-\rh);
    \node[font=\scriptsize] at (0.45,-\i*\rh-0.5*\rh) {$\lbl$};
  }
  % the residual entry g_{j+1}
  \draw[simRust, very thick] (0,-4*\rh) rectangle ++(0.9,-\rh);
  \fill[simBgRust] (0,-4*\rh) rectangle ++(0.9,-\rh);
  \draw[simRust, very thick] (0,-4*\rh) rectangle ++(0.9,-\rh);
  \node[font=\scriptsize, text=simRust] at (0.45,-4*\rh-0.5*\rh) {$g_{j+1}$};
  % brace + label
  \draw[decorate, decoration={brace, amplitude=4pt}, simBlue]
    (1.05,-0) -- (1.05,-4*\rh) node[midway, right=3pt, font=\scriptsize, text=simBlue, align=left]
    {fitted part:\\driven to $0$\\by back-solve};
  \draw[decorate, decoration={brace, amplitude=4pt, mirror}, simRust]
    (-0.05,-4*\rh) -- (-0.05,-5*\rh);
  \node[font=\scriptsize, text=simRust, anchor=east, align=right] at (-0.15,-4.5*\rh)
    {the residual\\$\norm{\vr_m} = |g_{j+1}|$\\\emph{free!}};
  % the "no Ax needed" annotation
  \node[extern, minimum width=30mm, align=center] at (5.0,-1.0)
    {\textbf{not computed:}\\$\vx_m = \vx_0+\mat V_m\vy_m$\\$\mat A\vx_m$, \;$\vb - \mat A\vx_m$};
  \draw[-{Stealth[length=2mm]}, simGray, densely dotted] (4.0,-2.6) -- (1.4,-2.3);
  \node[font=\scriptsize\itshape, text=simGray] at (4.6,-2.8) {avoided every step};
\end{tikzpicture}
\caption[The residual is a free byproduct]{Reading the residual off the rotated
right-hand side. After the rotation stream acts on $\vg = \mat Q_m(\beta\vec e_1)$, the
upper $m$ entries are the fittable part---a back-solve can drive $\norm{\vg_{0:m} -
\mat R_m\vy}$ to zero---and the single last entry $g_{j+1}$ (rust) is the part no
choice of $\vy$ can fit. Its magnitude \emph{is} the GMRES residual norm, by
\cref{eq:gv-split}. The expensive alternative on the right---forming the iterate
$\vx_m$, multiplying by $\mat A$, and subtracting from $\vb$---is never needed at all.
The residual is a unitary byproduct of triangularizing the matrix.}
\label{fig:gv-residual}
\end{figure}

\begin{keyidea}
The least-squares residual is a \emph{free byproduct} of the running QR. Each plane
rotation is orthogonal, so applying the rotation stream to the right-hand side $\beta\,
\vec e_1$ preserves its norm and concentrates the entire un-fittable component into one
entry. After the right-hand-side update, the magnitude of that entry, $\abs{g_{j+1}}$,
\emph{is} the GMRES residual norm $\norm{\beta\,\vec e_1 - \bar{\mat H}_m\vy_m}$---no
iterate $\vx_m$ formed, no matrix--vector product $\mat A\vx_m$ computed, no residual
evaluated. This is why GMRES can monitor convergence at essentially zero cost: the
quantity the convergence test needs falls out of the triangularization the solver was
already doing.
\end{keyidea}

\section{The back-solve: the terminal consumer}
\label{sec:gv-backsolve}

The stream runs every Arnoldi step, cheaply, and tells us when to stop. But it never
forms the answer. When we finally want the actual coordinate vector $\vy_m$---at
convergence, or at a restart boundary where GMRES discards the basis and begins
again---we run the \emph{back-solve} once. The triangularized matrix has made this easy:
$\mat R_m$ (the leading $m\times m$ block of $\bar{\mat R}_m$) is upper-triangular, so
$\mat R_m\vy = \vg_{0:m}$ is solved by back-substitution, bottom row up:
\begin{equation}
\label{eq:gv-backsub}
  y_i = \frac{1}{R_{ii}}\!\left(g_i - \sum_{k>i} R_{ik}\,y_k\right),
  \qquad i = m{-}1, m{-}2, \dots, 0.
\end{equation}
Each $y_i$ uses only the $y_k$ already computed below it; the last unknown $y_{m-1} =
g_{m-1}/R_{m-1,m-1}$ comes first and the rest follow upward. This is an $O(m^2)$ scalar
solve on the small dense triangle---trivial beside one field-sized matrix--vector
product. Once $\vy_m$ is in hand, the iterate correction is reconstructed against the
Arnoldi basis,
\begin{equation}
\label{eq:gv-reconstruct}
  \vx_m = \vx_0 + \mat V_m\,\vy_m = \vx_0 + \sum_{k=0}^{m-1} y_k\,\vv_k,
\end{equation}
a single \code{linear\_combination} (\cref{ch:combinators}) over the basis vectors. For
flexible GMRES (FGMRES), the only change is the basis: the correction is reconstructed
against the preconditioned basis $\mat Z_m$ rather than $\mat V_m$, so
$\vx_m = \vx_0 + \mat Z_m\vy_m$. The running-QR stream, the rotations, and the
back-solve are bit-for-bit identical between GMRES and FGMRES; only the basis the final
sum reads differs.

\begin{figure}[t]
\centering
\begin{tikzpicture}[font=\footnotesize]
  % upper-triangular R with back-substitution sweep arrow
  \def\cw{0.66} \def\rh{0.56}
  \foreach \c in {0,1,2,3}{\foreach \r in {0,...,3}{
      \draw[simGray!30] (\c*\cw,-\r*\rh) rectangle ++(\cw,-\rh);}}
  % shade upper triangle
  \foreach \c/\minr in {0/0,1/0,2/0,3/0}{
    \foreach \r in {0,...,\c}{
      \fill[simBgBlue] (\c*\cw,-\r*\rh) rectangle ++(\cw,-\rh);}}
  \foreach \c in {0,1,2,3}{\foreach \r in {0,...,3}{
      \draw[simGray!30] (\c*\cw,-\r*\rh) rectangle ++(\cw,-\rh);}}
  % diagonal entries
  \foreach \i in {0,1,2,3}{\node[font=\scriptsize] at (\i*\cw+0.5*\cw,-\i*\rh-0.5*\rh) {$R_{\i\i}$};}
  \node[font=\scriptsize, text=simGray] at (2.2*\cw,0.30) {$\mat R_m$ (upper-tri)};
  % RHS vector g
  \node[font=\scriptsize] at (4.7*\cw,0.30) {$\vg_{0:m}$};
  \foreach \r in {0,1,2,3}{
    \draw[simGray!45] (4.4*\cw,-\r*\rh) rectangle ++(0.8*\cw,-\rh);
    \node[font=\scriptsize] at (4.8*\cw,-\r*\rh-0.5*\rh) {$g_{\r}$};}
  % y vector
  \node[font=\scriptsize] at (6.2*\cw,0.30) {$\vy$};
  \foreach \r in {0,1,2,3}{
    \draw[simGreen!50] (5.9*\cw,-\r*\rh) rectangle ++(0.8*\cw,-\rh);
    \node[font=\scriptsize, text=simGreen] at (6.3*\cw,-\r*\rh-0.5*\rh) {$y_{\r}$};}
  % the bottom-up sweep arrow
  \draw[-{Stealth[length=2.4mm]}, simRust, very thick] (6.3*\cw,-3.7*\rh) -- (6.3*\cw,-0.2*\rh);
  \node[font=\scriptsize, text=simRust, anchor=west, align=left] at (7.0*\cw,-2*\rh)
    {back-substitute\\\emph{bottom row up}:\\$y_{m-1}$ first,\\then sweep up};
\end{tikzpicture}
\caption[The back-solve as the terminal consumer]{The back-solve, run once at
convergence or restart. The triangularized $\mat R_m$ (left, upper-triangular) and the
fitted part $\vg_{0:m}$ of the rotated right-hand side feed a back-substitution
(\cref{eq:gv-backsub}) that sweeps from the bottom row upward, computing $y_{m-1}$
first and each higher $y_i$ from the ones already found. The result $\vy$ (green)
reconstructs the iterate $\vx_m = \vx_0 + \mat V_m\vy$ via one
\code{linear\_combination}. This is the stream's terminal \emph{consumer}: it runs
once per restart cycle, not once per step.}
\label{fig:gv-backsolve}
\end{figure}

\subsection{Producer and consumer: two cadences}
\label{sec:gv-producer-consumer}

The chapter's architecture is a clean split between a streaming \emph{producer} and a
terminal \emph{consumer}, and the two run at different cadences. The producer is the
running-QR stream of \cref{sec:gv-stream}: it fires \emph{every Arnoldi step}, consuming
one column, appending one rotation, and emitting one residual scalar. The consumer is the
back-solve of \cref{sec:gv-backsolve}: it fires \emph{once per restart cycle}, at the
moment we actually want the answer. \Cref{fig:gv-prodcons} draws the two cadences
side by side.

This split is not incidental; it is why the running QR is worth maintaining. If we
solved the least-squares problem from scratch each step, every step would pay the
back-solve cost \emph{and} the re-factorization cost. By streaming the factorization and
deferring the back-solve, each step pays only the cheap incremental rotation
work---which also happens to hand us the residual for free---and the one back-solve is
amortized over the whole restart cycle. The producer monitors; the consumer materializes;
and the convergence test that decides when to call the consumer reads the producer's free
byproduct.

\begin{figure}[t]
\centering
\begin{tikzpicture}[font=\footnotesize]
  % producer: a chain of per-step stream updates, each emitting beta
  \node[font=\small\bfseries, text=simGreen] at (3.0,2.3) {producer: the stream (every step)};
  \foreach \j/\x in {0/0,1/2.0,2/4.0,3/6.0}{
    \node[opbox, minimum width=15mm, fill=simBgGreen, draw=simGreen] (u\j) at (\x,1.3)
      {step \j};
    \node[font=\scriptsize, text=simRust] at (\x,0.55) {$|g_{j{+}1}|$};
    \draw[-{Stealth[length=1.6mm]}, simRust] (\x,1.0) -- (\x,0.75);
  }
  \foreach \j/\jn in {0/1,1/2,2/3}{
    \draw[flow, simGreen] (u\j.east) -- (u\jn.west);}
  \node[font=\scriptsize\itshape, text=simGreen, anchor=west] at (6.9,1.3) {$\dots$};
  % the convergence test reading the byproduct
  \node[font=\scriptsize\itshape, text=simRust, align=center] at (3.0,0.05)
    {convergence test reads $|g_{j+1}|$ at every step (free)};
  % consumer: the back-solve, fired once
  \node[font=\small\bfseries, text=simBlue] at (3.0,-1.2) {consumer: the back-solve (once at restart)};
  \node[product, minimum width=24mm] (bs) at (3.0,-2.1)
    {back-solve $\to \vy_m \to \vx_m$};
  \draw[-{Stealth[length=2mm]}, simBlue, densely dashed] (u3.south) .. controls +(0,-1.3) and +(1.5,1.0) .. (bs.east)
    node[midway, right, font=\scriptsize, text=simBlue] {fire \emph{once}};
\end{tikzpicture}
\caption[Producer versus consumer: two cadences]{The two cadences of the least-squares
core. The \emph{producer} (top, green) is the running-QR stream: it runs once per
Arnoldi step, threading the rotation list forward and emitting the residual $|g_{j+1}|$
each step, which the convergence test reads for free. The \emph{consumer} (bottom, blue)
is the back-solve: it fires \emph{once}, at convergence or a restart boundary, to
materialize $\vy_m$ and reconstruct the iterate $\vx_m$. Streaming the factorization and
deferring the back-solve is the whole reason the running QR beats re-solving from
scratch: the per-step cost is cheap, and the one expensive consumer is amortized.}
\label{fig:gv-prodcons}
\end{figure}

\section{The whole inner solve, worked}
\label{sec:gv-worked}

We now run the entire machine on a concrete small problem: two stream steps, the
residual read off for free, and a back-solve for the answer. This is the GMRES inner
solve in miniature.

\begin{workedexample}{Streaming QR on a $3\times2$ Hessenberg, with the free residual}{gv-stream}
Take $\beta = 5$ and the $3\times2$ upper-Hessenberg matrix that Arnoldi might have
produced after two steps:
\[
  \bar{\mat H}_2 = \begin{pmatrix} 3 & 1 \\ 4 & 5 \\ 0 & 12 \end{pmatrix},
  \qquad \vg^{(0)} = \beta\,\vec e_1 = \begin{pmatrix} 5 \\ 0 \\ 0 \end{pmatrix}.
\]
We want $\min_\vy\norm{\beta\vec e_1 - \bar{\mat H}_2\vy}$, and we want the residual for
free along the way.

\medskip\emph{Step 0 --- the first column} $\vh_0 = (3,4,0)^{\!\top}$.
There are no stored rotations yet, so the \emph{replay} sub-step is empty. \emph{Generate}
from the pair $(h_0, h_1) = (3,4)$: this is exactly \cref{wex:gv-generate}, giving
$(c_0, s_0) = (0.6, 0.8)$ and $r = 5$. \emph{Apply} to the column zeros the
sub-diagonal: $(3,4)\mapsto(5,0)$, so column 0 of $\mat R$ is $R_{00} = 5$. \emph{Apply
to the RHS} pair $(g_0, g_1) = (5,0)$:
\[
  \begin{pmatrix} 0.6 & 0.8 \\ -0.8 & 0.6 \end{pmatrix}\!\begin{pmatrix} 5 \\ 0 \end{pmatrix}
  = \begin{pmatrix} 3 \\ -4 \end{pmatrix},
  \quad\text{so}\quad \vg^{(1)} = (3,\,-4,\,0)^{\!\top}.
\]
The residual after one step is $\abs{g_1} = \abs{-4} = 4$. (Not yet converged; carry on.)

\medskip\emph{Step 1 --- the second column} $\vh_1 = (1,5,12)^{\!\top}$.
Now there \emph{is} a stored rotation, so we \emph{replay} $(c_0,s_0) = (0.6,0.8)$ on the
top pair $(h_0, h_1) = (1,5)$:
\[
  \begin{pmatrix} 0.6 & 0.8 \\ -0.8 & 0.6 \end{pmatrix}\!\begin{pmatrix} 1 \\ 5 \end{pmatrix}
  = \begin{pmatrix} 0.6 + 4 \\ -0.8 + 3 \end{pmatrix}
  = \begin{pmatrix} 4.6 \\ 2.2 \end{pmatrix},
\]
leaving the replayed column $(4.6,\,2.2,\,12)^{\!\top}$ (the bottom entry $12$ is
untouched by the replay---rotation 0 only mixes rows 0 and 1). The off-diagonal
$R_{01} = 4.6$ is now fixed. \emph{Generate} from the new sub-diagonal pair
$(h_1, h_2) = (2.2, 12)$:
\[
\begin{aligned}
  r &= \sqrt{2.2^2 + 12^2} = \sqrt{4.84 + 144} = \sqrt{148.84} = 12.2,\\
  c_1 &= \frac{2.2}{12.2} \approx 0.1803,\qquad s_1 = \frac{12}{12.2} \approx 0.9836.
\end{aligned}
\]
\emph{Apply} to the column zeros the sub-diagonal: $(2.2,12)\mapsto(12.2,0)$, so
$R_{11} = 12.2$. \emph{Apply to the RHS} pair $(g_1, g_2) = (-4, 0)$:
\[
  \begin{pmatrix} c_1 & s_1 \\ -s_1 & c_1 \end{pmatrix}\!\begin{pmatrix} -4 \\ 0 \end{pmatrix}
  = \begin{pmatrix} 0.1803\cdot(-4) \\ -0.9836\cdot(-4) \end{pmatrix}
  \approx \begin{pmatrix} -0.721 \\ 3.934 \end{pmatrix},
\]
so $\vg^{(2)} \approx (3,\,-0.721,\,3.934)^{\!\top}$.

\medskip\emph{Read the residual, free.} The new last entry is
$\abs{g_2} \approx 3.934$. By \cref{eq:gv-free} this \emph{is} the least-squares
residual norm $\min_\vy\norm{\beta\vec e_1 - \bar{\mat H}_2\vy}$---and we have computed
no iterate and no $\mat A\vx$. (Cross-check: the rotations are unitary, so the norm of
$\vg$ is preserved: $\norm{\vg^{(2)}}^2 \approx 3^2 + 0.721^2 + 3.934^2 \approx 9 +
0.520 + 15.48 = 25.0 = \beta^2$. The length is conserved, as it must be.)

\medskip\emph{Back-solve for} $\vy$. The fitted part is the $2\times2$ triangular system
\[
  \mat R_2 = \begin{pmatrix} 5 & 4.6 \\ 0 & 12.2 \end{pmatrix},
  \qquad \vg_{0:2} = \begin{pmatrix} 3 \\ -0.721 \end{pmatrix}.
\]
Back-substitute bottom-up (\cref{eq:gv-backsub}). The last unknown first:
\[
  y_1 = \frac{g_1}{R_{11}} = \frac{-0.721}{12.2} \approx -0.0591.
\]
Then sweep up:
\[
  y_0 = \frac{1}{R_{00}}\big(g_0 - R_{01}\,y_1\big)
      = \frac{1}{5}\big(3 - 4.6\cdot(-0.0591)\big)
      = \frac{1}{5}\big(3 + 0.272\big) \approx 0.654.
\]
So $\vy_2 \approx (0.654,\,-0.0591)^{\!\top}$, and the GMRES iterate is $\vx_2 = \vx_0 +
y_0\,\vv_1 + y_1\,\vv_2$---one \code{linear\_combination} over the two Arnoldi basis
vectors. The residual we read off for free, $\approx 3.934$, is the norm we would
measure if we formed that $\vx_2$ and computed $\norm{\vb - \mat A\vx_2}$---but we never
had to.
\end{workedexample}

\Cref{wex:gv-stream} is the entire GMRES inner solve at $m=2$: two stream steps building
$\mat R$ one column at a time, the residual read off the rotated right-hand side without
any field-sized work, and a single back-solve producing the coordinates. Scale $m$ from
$2$ to $40$ and nothing changes but the size of the small triangle; the cost of the
field-sized operations---the Arnoldi matrix--vector products of \cref{ch:orthog}---lives
entirely outside this chapter.

\section{Where this sits in the machine}
\label{sec:gv-context}

It is worth placing the running QR in the larger GMRES picture one last time. The Arnoldi
process of \cref{ch:orthog} is the \emph{expensive} part: each step does a field-sized
matrix--vector product $\mat A\vv_k$ and orthogonalizes against the basis, both $O(N)$ or
worse. The least-squares core of this chapter is the \emph{cheap} part: it works only on
the small dense Hessenberg and the small right-hand side, $O(m)$ scalars per step. The two
interlock once per iteration---Arnoldi hands the running QR a new column, the running QR
hands back a residual norm that decides whether Arnoldi runs again---and at the end the
back-solve crosses back into field space exactly once, to reconstruct $\vx_m$ from the
basis. The producer monitors cheaply; the consumer materializes rarely; and the residual,
the one quantity GMRES most needs and would otherwise pay dearly for, falls out of the
factorization for nothing. That free residual is the elegant heart of the method, and it
is why GMRES, despite its appetite for memory, remains the workhorse of \cref{ch:gmres}.

\summarysection
GMRES needs, at each step $m$, the solution of a small upper-Hessenberg least-squares
problem $\min_\vy\norm{\beta\,\vec e_1 - \bar{\mat H}_m\vy}$. Rather than re-factor from
scratch, the solver maintains a \emph{running QR} factorization, updated one column at a
time by a stream of $2\times2$ Givens (plane) rotations. A Givens rotation $G(c,s)$ with
$c = a/r$, $s = b/r$, $r = \sqrt{a^2+b^2}$ zeros one sub-diagonal entry; production code
computes $r$ by the scaled, overflow-safe formula. The stream processes each arriving
column in four sub-steps---\emph{replay} the stored rotations, \emph{generate} a new one,
\emph{apply} it to the column, \emph{apply} it to the right-hand side---and is exactly a
\code{foldl} over columns, with the rotation list as the growing carry; the replay must
precede the generate, because plane rotations do not commute. The decisive payoff: after
the right-hand-side update, the magnitude of its new last entry $\abs{g_{j+1}}$ \emph{is}
the GMRES residual norm, a free byproduct of unitarity---no iterate, no $\mat A\vx$, no
residual evaluation. The back-solve, run once per restart cycle, solves the small
triangular system by back-substitution and reconstructs $\vx_m = \vx_0 + \mat V_m\vy$
(or $\mat Z_m\vy$ for FGMRES) by one \code{linear\_combination}. The stream is a
streaming producer running every step; the back-solve is a terminal consumer running
once.

\furtherreading
The running-QR-via-Givens implementation of GMRES is treated in
Saad~\citep{saad2003}, \S6.5.3, which is the closest reference to the stream as Palace
realizes it, including the free residual. The Givens rotation itself, its stable
generation, and its use in QR factorization are in Golub \& Van
Loan~\citep{golubvanloan2013}, \S5.1 (rotations) and \S5.2 (QR); the overflow-safe scaled
generation follows the LAPACK convention discussed there. Trefethen \&
Bau~\citep{trefethenbau1997}, Lectures 10 and 35, give an unusually clear geometric
account of orthogonal triangularization and of GMRES as residual minimization over a
Krylov subspace. The complex-arithmetic conventions and the precise scaled formula are
collected in \cref{app:numerics}.

\exercisessection
\begin{enumerate}[nosep]
  \item \emph{Generate by hand.} Compute the Givens rotation $(c,s)$ that zeros the
  second component of $(a,b) = (5,12)$. Give $r$, $c$, $s$, verify $c^2+s^2=1$, and check
  that applying $G(c,s)$ to $(5,12)$ yields $(13,0)$.
  \item \emph{Orthogonality preserves length.} Show directly from \cref{eq:gv-matrix}
  that $G(c,s)$ is orthogonal ($G^{\!\top}G = I$) whenever $c^2+s^2=1$, and conclude that
  $\norm{G\vu} = \norm{\vu}$ for every 2-vector $\vu$. Where in the chapter is this fact
  the reason the residual is free?
  \item \emph{Why scale.} In double precision the overflow threshold is about
  $1.8\times10^{308}$. Take $a = b = 10^{200}$. Show that the literal $r=\sqrt{a^2+b^2}$
  overflows (compute $a^2$), while the scaled formula \eqref{eq:gv-scaled} returns the
  correct finite $r = 10^{200}\sqrt2$. By how large a factor does $a^2$ exceed the
  threshold?
  \item \emph{Replay order matters.} In \cref{wex:gv-stream}, redo step 1 but
  \emph{skip} the replay of $(c_0,s_0)$ and generate the new rotation directly from the
  raw column $(1,5,12)$. Show that the resulting $\mat R$ is no longer the QR factor of
  $\bar{\mat H}_2$ (e.g.\ the residual you would read off no longer matches the true
  least-squares residual). Explain in one sentence which load-bearing invariant you
  violated.
  \item \emph{Read the residual.} A GMRES run reports, after its rotation stream at step
  $m$, a rotated right-hand side whose last entry is $g_{m} = -1.5\times10^{-7}$ with
  $\beta = 2.0$. What is the current relative residual $\norm{\vr_m}/\norm{\vb}$
  (assume $\vx_0=0$ so $\beta=\norm{\vb}$)? Would a tolerance of $10^{-6}$ declare
  convergence? Note that you answered this without forming $\vx_m$.
  \item \emph{Back-solve.} Given $\mat R_2 = \begin{psmallmatrix} 2 & 1 \\ 0 & 4
  \end{psmallmatrix}$ and the fitted right-hand side $\vg_{0:2} = (6, 8)^{\!\top}$, solve
  $\mat R_2\vy = \vg_{0:2}$ by the back-substitution of \cref{eq:gv-backsub}. Which
  unknown do you compute first, and why?
  \item \emph{The fold made explicit.} Write \code{ls\_update\_column} as a fold step
  $st' = \code{step}(st, \vh)$ and identify, for the running QR, (a) the carry $st$,
  (b) the per-element input $\vh$, and (c) the part of the carry that \emph{grows} with
  each step. Why does the presence of the replay sub-step force this to be a
  \code{foldl} (a chain) rather than a \code{map} (an independent fan-out)? Refer to the
  map/fold split of \cref{ch:combinators}.
  \item \emph{Producer versus consumer.} The stream runs every Arnoldi step and the
  back-solve runs once per restart cycle. If a GMRES run takes $40$ steps before a
  restart, how many times does each fire? Explain why moving the back-solve inside the
  per-step loop (so it fires every step) would be correct but wasteful, and estimate the
  factor of extra back-solve work it would incur.
\end{enumerate}
